\begin{flushright}
\textit{Because all the data values lie within the bounds on the normal probability plot (see the discussion in Section 3.2.4), we can assume that the data have approximate normality. We have $\bar{X} = 285.7$ and $s = 43.9$. Also, $n = 12$. For $\alpha = 0.05$, $t_{0.025, 11} = 2.201$. A 95\% CI for $\mu$ is,}
\end{flushright}

\[
\bar{x} \pm t_{(\alpha/2,n-1)} \frac{s}{\sqrt{n}} = 285.7 \pm 2.201 \left( \frac{43.9}{\sqrt{12}} \right)
\]

\textit{Hence, a 95\% CI for $\mu$, the average ozone level over the Earth, lies in (257.81, 313.59).}

\section*{EXERCISES 5.5}

\begin{enumerate}
    \item[\textbf{5.5.1.}] A survey indicates that it is important to pay attention to truth in political advertising. Based on a survey of 1200 people, 35\% indicated that they found political advertisements to be untrue; 60\% say that they will not vote for candidates whose advertisements are judged to be untrue; and of this latter group, only 15\% ever complained to the media or to the candidate about their dissatisfaction.
    \begin{enumerate}
        \item[(a)] Find a 95\% CI for the percentage of people who find political advertising to be untrue.
        \item[(b)] Find a 95\% CI for the percentage of voters who will not vote for candidates whose advertisements are considered to be untrue.
        \item[(c)] Find a 95\% CI for the percentage of those who avoid voting for candidates whose advertisements are considered untrue and who have complained to the media or to the candidate about the falsehoods in commercials.
        \item[(d)] For each case above, interpret the results and state any assumptions you have made.
    \end{enumerate}

    \item[\textbf{5.5.2.}] Many mutual funds use an investment approach involving owning stocks whose price/earnings multiples (P/Es) are less than the P/E of the S\&P 500. The following data give P/Es of 49 companies that a randomly selected mutual fund owns in a particular year.

    \begin{center}
    \begin{tabular}{ccccccccccc}
    6.8 & 5.6 & 8.5 & 8.5 & 8.4 & 7.5 & 9.3 & 9.4 & 7.8 & 7.1 \\
    9.9 & 9.6 & 9.0 & 9.4 & 13.7 & 16.6 & 9.1 & 10.1 & 10.6 & 11.1 \\
    8.9 & 11.7 & 12.8 & 11.5 & 12.0 & 10.6 & 11.1 & 6.4 & 12.3 & 12.3 \\
    11.4 & 9.9 & 14.3 & 11.5 & 11.8 & 13.3 & 12.8 & 13.7 & 13.9 & 12.9 \\
    14.2 & 14.0 & 15.5 & 16.9 & 18.0 & 17.9 & 21.8 & 18.4 & 34.3 & \\
    \end{tabular}
    \end{center}

    Find a 98\% CI for the mean P/E multiples. Interpret the result and state any assumptions you have made.

    \item[\textbf{5.5.3.}] Let $X_1, \ldots, X_n$ be a random sample from an $N(\mu, \sigma^2)$ distribution, with $\sigma^2$ known.
    \begin{enumerate}
        \item[(a)] Show that $\hat{\mu} = \bar{X}$ is an MLE of the population mean $\mu$.
        \item[(b)] Show that:
        \[
        P\left( \bar{X} - \frac{2\sigma}{\sqrt{n}} < \mu < \bar{X} + \frac{2\sigma}{\sqrt{n}} \right) = 0.954.
        \]
        \item[(c)] Let
        \[
        P\left( \bar{X} - \frac{k\sigma}{\sqrt{n}} < \mu < \bar{X} + \frac{k\sigma}{\sqrt{n}} \right) = 0.90.
        \]
        Find $k$.
    \end{enumerate}

    \item[\textbf{5.5.4.}] Let the observed mean of a sample of size 45 be $\bar{x} = 68.51$ from a distribution having variance 110. Find a 95\% CI for the true mean $\mu$ and interpret the result and state any assumptions you have made.

    \item[\textbf{5.5.5.}] In a random sample of 50 college seniors, 18 indicated that they were planning to pursue a graduate degree. Find a 98\% CI for the true proportion of all college seniors planning to pursue a graduate degree, and interpret the result, and state any assumptions you have made.

    \item[\textbf{5.5.6.}] DVD players coming off an assembly line are automatically checked to make sure they are not defective. The manufacturer wants an interval estimate of the percentage of DVD players that fail the testing procedure. Compute a 90\% CI, based on a random sample of size 105 in which 17 DVD players failed the testing procedure. Also, interpret the result and state any assumptions you have made.
\end{enumerate}
